\subsection{The first complex norm transition at seventeen}

\begin{statement}{E.57}{Proved}{Actual class-group corrections and the complete first complex norm map}
For the first cocycle $3\kappa_3(17)$ in E.46, the
first two positional algebras and their cubing map are
\[
 \begin{gathered}
 R_1=\mathbb Q[X]/(X^3-17^3)=\mathbb Q\times K,
                   \qquad K=\mathbb Q(\zeta),\quad\zeta^2+\zeta+1=0,\\
 R_2=\mathbb Q[X]/(X^9-17^3)=F\times L,\\
 F=\mathbb Q(\alpha),\quad\alpha^3=17,\qquad
 L=K(\beta),\quad\beta^3=17\zeta,\\
 R_1\longrightarrow R_2,\quad X\longmapsto X^3,\qquad
 N_{R_2/R_1}=(N_{F/\mathbb Q},N_{L/K}).
 \end{gathered}
 \tag{E57.1}
\]
In particular, $L=\mathbb Q(\beta)$ has defining
polynomial $\beta^6+17\beta^3+289$.
The following formulas evaluate its entire norm
map, including the integral ideal-class corrections.
Every completion in this statement is
$A^{\wedge_3}=\varprojlim_n A/3^nA$; restricted
infinite coordinates carry that inverse-limit
topology.

Put $\lambda=1-\zeta$, $\pi=2-\zeta$,
$\bar\pi=3+\zeta$, and
\[
 \begin{gathered}
 \mathfrak p_1=(7,\beta-3),\quad
 \mathfrak p_2=(7,\beta+1),\quad
 \mathfrak p_3=(7,\beta+2),\\
 c=\beta+3\zeta+1,\qquad
 g_1=-2-3\zeta-\zeta\beta,\qquad
 g_2=-3-\zeta+\zeta\beta .
 \end{gathered}
 \tag{E57.2}
\]
The class group and its specified generators are
\[
 \operatorname{Cl}(L)
   =(\mathbb Z/3)[\mathfrak p_1]
           \oplus(\mathbb Z/3)[\mathfrak p_2],
 \qquad [\mathfrak p_3]=-[\mathfrak p_1]-[\mathfrak p_2].
 \tag{E57.3}
\]
There is one prime $\mathfrak q_3$ above three,
one $\mathfrak q_{17}$ above seventeen,
and one degree-three prime $\mathfrak q_{\bar\pi}$
above $\bar\pi$. The literal principal-ideal
identities are
\[
 \begin{gathered}
 (c)=\mathfrak q_3,\quad
 (\beta)=\mathfrak q_{17},\quad
 (\bar\pi)=\mathfrak q_{\bar\pi},\quad
 (\pi)=\mathfrak p_1\mathfrak p_2\mathfrak p_3,\\
 (g_1)=\mathfrak p_1^3,\qquad
 (g_2)=\mathfrak p_2^3 .
 \end{gathered}
 \tag{E57.4}
\]
Thus a divisor supported above $3,7,17$, with
coefficients $n_i$ at $\mathfrak p_i$, is principal
precisely when
$n_1\equiv n_2\equiv n_3\pmod3$.

Let $S_K$ consist of the primes above $3,7,17$,
and let $S_L$ be its inverse image. Two actual
fundamental units of $L$ are
\[
 \begin{aligned}
 e_1&=-18+3\beta+3\beta^2
              -18\beta^3/17-5\beta^4/17+2\beta^5/17,\\
 e_2&=-18-5\beta+2\beta^2
              -\beta^3/17-8\beta^4/17-\beta^5/17.
 \end{aligned}
 \tag{E57.5}
\]
The source $S$-unit completion has the specified
basis
\[
 \mathcal O_{L,S_L}^{\times,\wedge_3}
   =\langle e_1,e_2,c,\beta,\bar\pi,\pi,g_1,g_2\rangle_{\mathbb Z_3}
                                \oplus\langle\zeta\rangle_{\mathbb Z/3}.
 \tag{E57.6}
\]
The target basis is
$\lambda,17,\pi,\bar\pi$ over $\mathbb Z_3$ and
$\zeta$ of order three. Relative norms are exactly
\[
 \begin{array}{c|ccccccccc}
 v&e_1&e_2&c&\beta&\bar\pi&\pi&g_1&g_2&\zeta\\ \hline
 N(v)&1&1&\lambda&17\zeta&\bar\pi^3&\pi^3&
                   -\zeta^2\pi^3&\zeta^2\pi^3&1 .
 \end{array}
 \tag{E57.7}
\]
The sign is killed by the explicitly defined
three-primary completion. In these coordinates
the norm matrix has its last row read modulo three:
\[
 \begin{pmatrix}
 0&0&1&0&0&0&0&0&0\\
 0&0&0&1&0&0&0&0&0\\
 0&0&0&0&0&3&3&3&0\\
 0&0&0&0&3&0&0&0&0\\
 0&0&0&1&0&0&2&2&0
 \end{pmatrix}.
 \tag{E57.8}
\]
Consequently its full $S$-unit kernel and cokernel
are
\[
 \begin{aligned}
 \ker N&=
  \langle e_1,e_2,-g_1/g_2,(g_2/\pi)^3\rangle_{\mathbb Z_3}
                         \oplus\langle\zeta\rangle_{\mathbb Z/3}
       \simeq\mathbb Z_3^4\oplus\mathbb Z/3,\\
 \operatorname{coker}N&\simeq(\mathbb Z/3)^2,
           \qquad\text{coordinates }(v_\pi\bmod3,v_{\bar\pi}\bmod3).
 \end{aligned}
 \tag{E57.9}
\]
There is a concrete correction with zero relative
divisor norm:
\[
 w=(g_1/\pi)^2,\qquad
 N(w)=\zeta,\qquad
 \operatorname{div}(w)=
     4\mathfrak p_1-2\mathfrak p_2-2\mathfrak p_3 .
 \tag{E57.10}
\]
By contrast, the ordinary-unit norm image is
exactly $\{1,-1\}$. Formula (E57.10) specifies
the extra valuations which supply this root norm.

Here are full global coordinates, so that the
calculation also specifies every column outside
$S$. For every prime $\mathfrak P\notin S_L$,
compute its unique class coordinates
$a_{\mathfrak P},b_{\mathfrak P}\in\{0,1,2\}$
in (E57.3), and choose the actual principal
generator $q_{\mathfrak P}$ of
\[
 (q_{\mathfrak P})=
       \mathfrak P\mathfrak p_1^{-a_{\mathfrak P}}
                         \mathfrak p_2^{-b_{\mathfrak P}} .
 \tag{E57.11}
\]
Class logarithms and generators are computed by
the certified ideal algorithm retained with
this statement. Fix a generator
$p_{\mathfrak p}$ for each prime
$\mathfrak p\notin S_K$ in the Eisenstein
principal ideal domain. If $\mathfrak P$ lies
over $\mathfrak p$, let $f_{\mathfrak P}$ be its
relative residue degree. The exact element
norm determines $\varepsilon_{\mathfrak P}\in\mathbb Z/3$ by
\[
 N(q_{\mathfrak P})
       =\pm\zeta^{\varepsilon_{\mathfrak P}}\,
          p_{\mathfrak p}^{f_{\mathfrak P}}\,
          \pi^{-a_{\mathfrak P}-b_{\mathfrak P}} .
 \tag{E57.12}
\]
These are evaluated number-field operations,
with the defining field, basis and class
generators fixed above.

Every element of $L^{\times,\wedge_3}$ has
unique coordinates
\[
 e_1^{t_1}e_2^{t_2}c^\alpha\beta^{b_0}
       \bar\pi^\delta\pi^\gamma g_1^\eta g_2^\theta
       \zeta^e\prod_{\mathfrak P\notin S_L}
                                      q_{\mathfrak P}^{n_{\mathfrak P}},
 \quad n_{\mathfrak P}\longrightarrow0
                       \text{ modulo every }3^m .
 \tag{E57.13}
\]
The complete global norm is obtained from
(E57.8) by the following explicit additions:
\[
 \begin{aligned}
 v_\lambda&=\alpha,& v_{17}&=b_0,\\
 v_\pi&=3(\gamma+\eta+\theta)
             -\sum_{\mathfrak P}(a_{\mathfrak P}+b_{\mathfrak P})
                                                   n_{\mathfrak P},
 &v_{\bar\pi}&=3\delta,\\
 v_{\mathfrak p}&=\sum_{\mathfrak P\mid\mathfrak p}
                                 f_{\mathfrak P}n_{\mathfrak P},&
 e_{\rm root}&=b_0+2\eta+2\theta+
                \sum_{\mathfrak P}\varepsilon_{\mathfrak P}
                                    n_{\mathfrak P}\pmod3 .
 \end{aligned}
 \tag{E57.14}
\]
All sums have their stated restricted convergence.
This is the actual first complex transition of
E.46. It imposes no conclusion about compatibility
with every subsequent norm transition, or about
lifting through the intrinsic kernels of (E47.8).
The manuscript anchor is \S51.16, folios
733v--736v, scans~172--178, formulas~(371)--(379),
Part~E(c),(d).
\end{statement}

\begin{proof}
Factoring the two polynomials in (E57.1) and
substituting $X^3$ gives all four stated
algebra maps. On the complex factor the
substitution is literally $\beta^3=17\zeta$,
so the associated corestriction in E.46 is
$N_{L/K}$ with this precise embedding.

The retained PARI/GP computation uses the
officially specified unconditional
\texttt{bnfcertify} procedure
\cite{E_PARICertifiedUnits}. It returns one
for both $K$ and $L$. It certifies the complete
class group $(3,3)$, the maximal-order basis
\[
 1,\ \zeta+1,\ \beta,\ \zeta\beta,\ \beta^2,\ \zeta\beta^2,
 \tag{E57.15}
\]
the two complete fundamental units in
(E57.5), and the torsion group $\mu_6$.
The class logarithms of the displayed
$\mathfrak p_1,\mathfrak p_2$ form an
invertible two-by-two matrix modulo three
on the certified class basis. Changing
by that actual matrix gives coordinates
$(1,0),(0,1),(2,2)$ for the three primes.
This proves (E57.3) with its specified
generators, not just its abstract order.

All remaining displayed ideal and element
equalities are checked independently of
PARI's class-group arithmetic. The Python
certificate constructs the exact
multiplication matrices in the integral
basis (E57.15), reduces modulo
$\beta^6+17\beta^3+289$, and compares their
Hermite lattices. It verifies (E57.4) and
the divisor in (E57.10), including its
negative exponents. Each relative norm is
also the independently computed determinant
\[
 N_{L/K}(a+b\beta+d\beta^2)
    =a^3+(17\zeta)b^3+(17\zeta)^2d^3
                                      -3(17\zeta)abd .
 \tag{E57.16}
\]
This gives every entry in (E57.7).
The official specification of
\texttt{bnfisprincipal} supplies both the
class logarithm and the principal
generator, and the code verifies the
resulting ideals and norms exactly.

For divisors supported on $S_L$, (E57.3)
reduces principality to
$n_1-n_3=n_2-n_3=0$ modulo three.
The lattice of these triples has basis
$(1,1,1),(3,0,0),(0,3,0)$, of index nine
in $\mathbb Z^3$. The other three prime
coordinates are independently principal
by (E57.4). Combining these six divisor
generators with the two fundamental units
and the torsion generator proves the
full basis (E57.6). No missing unit or
class index is absorbed into this basis.
The analogous target basis follows
from the Eisenstein principal ideal
domain and its units $\mu_6$.

Writing the source exponents in the order
of (E57.6), formula (E57.7) gives (E57.8).
For a norm-one vector its first, second,
and fourth target coordinates force
$\alpha=b_0=\delta=0$. The third says
$\gamma+\eta+\theta=0$, and its root
coordinate says $\eta+\theta\in3\mathbb Z_3$.
Write $\eta+\theta=3t$ and then
$\theta=3t-\eta$, $\gamma=-3t$.
These are exactly the four free
kernel generators in (E57.9), together
with the order-three root.
The root coordinate is attained by
(E57.10); the two remaining obstructions
are precisely the valuations at
$\pi,\bar\pi$ modulo three. Equivalently,
adjoining the target relation
$3e_{\rm root}=0$ to the integer norm
matrix gives Smith entries
$1,1,1,3,3$, as independently verified.
For ordinary units the certified basis
has norms $1,1$, while $N(-1)=-1$
and $N(\zeta)=1$. This proves the
exact ordinary-unit assertion and
the actual comparison (E57.10).

For the global formula, the primes in
$S_L$ generate the full class group.
Thus (E57.11) exists for every outside
prime. The resulting divisor sequence
has the explicit splitting
\[
 0\longrightarrow\mathcal O_{L,S_L}^{\times}
   \longrightarrow L^\times
   \xrightarrow{(v_{\mathfrak P})}
      \bigoplus_{\mathfrak P\notin S_L}\mathbb Z
   \longrightarrow0,\qquad
                1_{\mathfrak P}\longmapsto q_{\mathfrak P}.
 \tag{E57.17}
\]
This proves (E57.13), first as a
literal group decomposition and then
after completion for multiplication
by $3^n$. The completion of its direct
sum is exactly the restricted sequence
group stated there. This argument does
not use compactness of those infinite
quotients.

Taking the ideal norm of (E57.11)
gives
$\mathfrak p^{f_{\mathfrak P}}
(\pi)^{-a_{\mathfrak P}-b_{\mathfrak P}}$.
The difference between its two
displayed generators is a unit of
$K$, hence lies in $\mu_6$. This
proves the exact and computable
formula (E57.12). Taking its
coefficients gives (E57.14).
Since each class coordinate is in
$\{0,1,2\}$, the sums contributing
to the $S$ coordinates converge
for every restricted input.
There are at most three primes
over each target prime, and their
residue degrees are at most three,
so the outside target coordinates
again form a restricted family.
This proves continuity and the
entire asserted norm map.

The retained computation evaluates
all 31 outside-prime columns above
rational primes at most 61 other
than $3,7,17$. Each corrected
principal divisor, its exact
relative norm, and its root
coordinate is retained. The
Python certificate independently
checks these columns, 41 cubic
multiplication determinants,
nine ideal-lattice identities,
and the full Smith calculation.
The formulas (E57.11)--(E57.17)
prove the construction for every
prime; the finite computation
checks its concrete normalizations.
\end{proof}
